\documentclass[12pt]{article}
\usepackage[utf8]{inputenc}
\usepackage[T1]{fontenc}
\usepackage{hyperref}
\title{Quetzaltepec Mixe Language Resources for K-12 Curriculum Development}
\author{Compiled by Claude Sonnet 5.5 (claude-sonnet-5-5) for Andr\'as Kornai}
\date{}
\begin{document}
\maketitle

\section{Introduction}
Quetzaltepec Mixe (ISO 639-3 pxm; Glottocode quet1239) is a Mixe-Zoquean language of Oaxaca, Mexico. Search results describe it as threatened. The searches found very little material specific to this variety, and nothing that is a school curriculum resource. This report lists only what appeared in search results.

\section{Primary Educational Resources}
No primer, textbook or school programme specific to Quetzaltepec Mixe was found in the searches.

\section{Linguistic and Reference Materials}
\begin{itemize}
\item Swarthmore College linguistics course wiki page on Mixe. It notes that little modern literature exists on pxm beyond the work of Jany, and lists orthography, grammar discussion and Bible translations. \url{https://wikis.swarthmore.edu/ling073/Mixe}
\item Glottolog reference record for Bradley (1970), a Mixe-related linguistic reference. \url{https://glottolog.org/resource/reference/id/48857}
\item WALS reference record for Bradley (1970). \url{https://wals.info/refdb/record/Bradley-1970}
\item Paper on colonial Mixe texts with a brief grammar and glossary (Japanese with English sections); historical, not a modern Quetzaltepec grammar. \url{https://minpaku.repo.nii.ac.jp/record/4275/files/KH_016_2_003.pdf}
\end{itemize}

\section{Religious and Community Resources}
\begin{itemize}
\item Global Recordings Network page for Quetzaltepec Mixe (Bible-teaching audio). \url{https://globalrecordings.net/en/language/pxm}
\end{itemize}

\section{Digital and Audio Resources}
\begin{itemize}
\item Wikitongues language page for Quetzaltepec Mixe. \url{https://wikitongues.org/languages/pxm/}
\item Endangered Languages Project entry (notes on Quintana's colonial Confesionario covering Mixe). \url{https://www.endangeredlanguages.com/bn/node/112946}
\end{itemize}

\section{Curriculum Development Context}
Little documentation of Quetzaltepec Mixe was found. A curriculum would need to begin with community consultation and new documentation, drawing cautiously on materials for related Mixe varieties (for example Totontepec Mixe, which the Swarthmore page identifies as closest), which were not independently verified here.

\end{document}